\documentclass[11pt]{article}

% Useful packages
\usepackage[margin=1in]{geometry}
\usepackage{array}

\title{Investigacion Escalabilidad Horizontal y Vertical}
\author{Andreas Garoufalias}
\date{\today}

\begin{document}

\maketitle

\begin{abstract}
Este informe analiza las diferencias, características, ventajas y desventajas de la escalabilidad horizontal y vertical.
\end{abstract}

\section{Introduccion}

La escalabilidad es uno de los pilares mas grandes del mundo del software y hardware. Permitio la edad de la informacion - en el que vivimos hoy en dia - y dio introdujo la nube en nuestras vidas. Vamos a analizar los dos tipos de escalabilidad, horizontal y vertical, y entender la tecnologia que nos dio a todos una vida mas comoda.

\section{Escalabilidad Horizontal}

La escalabilidad horizontal, "scaling out," se refiere a agregar una o mas computadoras a un sistema para distribuir procesamiento. Las computadoras agregadas se les dicen "nodos," o servidores. Este tipo de escalabilidad se utiliza bastante en datacentres, datawarehouses, y otros tipos de centros de procesamiento de datos.
\newline

\textbf{Escalabilidad horizontal hace lo siguiente:}
\begin{itemize}
\item Agrega mas de un punto de fallo.
\item Aumenta las capacidades de un sistema casi infinitamente, con algunas excepciones.
\end{itemize}

\textbf{Este tipo de escalabilidad ofrece las siguientes caracteristicas:}
\begin{itemize}
\item Alta tolerancia a fallos
\item Capacidades aumentadas de almacenamiento y rendimiento.
\item Varios tipos de uso.
\item Complejidad alta de configuracion.
\end{itemize}

Finalmente, se utiliza mayormente por empresas medianas o grandes debido al costo de implementar este tipo de escalabilidad.

\section{Escalabilidad vertical}

La escalabilidad vertical, "scaling up," se refiere a aumentar las capacidades de una computadora, por ejemplo, cambiar el modelo del CPU por uno mas reciente y eficiente o agregar RAM ademas de el RAM que ya esta presente en la computadora. Este tipo de escalabilidad se utiliza mayormente por invididuales y empresas diminutas.
\newline

\textbf{Escalabilidad vertical hace lo siguiente:}
\begin{itemize}
	\item Agrega a o mejora las capacidades de una computadora individual.
	\item Permite una configuracion simple para realizarla la misma.
\end{itemize}

\textbf{La escalabilidad vertical ofrece las siguientes caracteristicas:}
\begin{itemize}
	\item Configuracion simple.
	\item Costo bajo comparado con la escalabilidad horizontal.
	\item Accesibilidad.
	\item Baja tolerancia a fallos debido a un solo punto de fallo.	
\end{itemize}


\section{Comparacion}


\begin{table}[h]
\centering
\begin{tabular}{p{8cm} p{8cm}}
\textbf{Horizontal} & \textbf{Vertical} \\[6pt]

Alta tolerancia a fallos. &
Baja tolerancia a fallos. \\

Complejidad alta de implementacion. &
Implementacion sencilla. \\

Altamente costoso en infraestructura. &
Costo relativamente bajo inicialmente. \\

Permite agregar multiples servidores o nodos. &
Solo permite mejorar un solo servidor. \\

Mayor escalabilidad a largo plazo. &
Escalabilidad limitada por el hardware. \\

Mayor complejidad en la gestion del sistema. &
Administracion mas simple del sistema. \\

Ideal para sistemas distribuidos y aplicaciones web grandes. &
Mas comun en sistemas pequenos o monoliticos. \\

\end{tabular}
\caption{Comparacion entre escalamiento horizontal y vertical}
\end{table}

\end{document}